% Chapter 4: System Design and Implementation
% Details the technical architecture, implementation, and design decisions
% for the thiLLMo OG-RAG system
\label{chap:design}

\section{System Architecture Overview}
\label{sec:system_architecture}

The thiLLMo (Ontology-Grounded RAG for Kikuyu) system implements the OG-RAG approach through a modular architecture comprising five primary subsystems: (1) Knowledge Graph Infrastructure, (2) Hybrid Retrieval Engine, (3) Context Builder, (4) LLM Integration Layer, and (5) Evaluation Pipeline. This chapter details the technical design and implementation of each component, the rationale for key architectural decisions, and the integration mechanisms that enable culturally faithful proverb translation.

\subsection{Design Principles}
\label{subsec:design_principles}

The system architecture adheres to four guiding principles:

\begin{description}
    \item[Modularity] Each subsystem operates as an independent module with well-defined interfaces, enabling isolated testing, incremental development, and potential reuse in other cultural NLP applications.
    
    \item[Transparency] All retrieval and generation steps produce interpretable outputs (e.g., which concepts matched, what cultural context was retrieved), facilitating error analysis and system debugging.
    
    \item[Scalability] The design accommodates future expansion from 100 proverbs to thousands, and from the wealth domain to broader cultural knowledge, without requiring architectural redesign.
    
    \item[Reproducibility] Configuration files, randomness seeds, and API parameters are versioned and logged, ensuring experimental results can be replicated.
\end{description}

\subsection{Technology Stack}
\label{subsec:tech_stack}

The system leverages the following core technologies:

\begin{itemize}
    \item \textbf{Neo4j 5.x (AuraDB)}: Cloud-hosted graph database for ontology storage and traversal queries
    \item \textbf{Python 3.11}: Primary implementation language
    \item \textbf{OpenAI API (GPT-4-turbo, GPT-4)}: LLM backend for translation generation
    \item \textbf{Google Gemini 2.0 Pro/Flash}: Alternative LLM for comparative evaluation and LLM-as-judge
    \item \textbf{Sentence-BERT}: Semantic embedding model for vector similarity search
    \item \textbf{NumPy/SciPy}: Statistical analysis and numerical computation
    \item \textbf{Matplotlib/Seaborn}: Visualization generation
\end{itemize}

The complete codebase is organized into four Python packages: \texttt{src/ontology} (schema management), \texttt{src/og-rag-system} (core translation pipeline), \texttt{src/evaluation} (metrics and baselines), and \texttt{src/data\_loading} (corpus ingestion).

\section{Knowledge Graph Infrastructure}
\label{sec:kg_infrastructure}

The knowledge graph serves as the structured semantic foundation for ontology-grounded retrieval, encoding the 847-concept ontology and 100-proverb corpus developed in Chapter~\ref{chap:methodology}.

\subsection{Neo4j Schema Design}
\label{subsec:neo4j_schema}

The graph schema implements a multi-layered structure optimizing for both ontological expressiveness and query efficiency.

\subsubsection{Node Types}
\label{subsubsec:node_types}

Four primary node types encode different knowledge facets:

\paragraph{Proverb Nodes} (\texttt{:Proverb}, 100 instances) represent individual proverbs with properties:
\begin{itemize}
    \item \texttt{proverb\_id}: Unique identifier (e.g., MW\_001)
    \item \texttt{kikuyu\_text}: Original Kikuyu proverb
    \item \texttt{expert\_translation}: Gold-standard English translation
    \item \texttt{expert\_cultural\_meaning}: Anthropological interpretation
    \item \texttt{expert\_business\_relevance}: Modern applicability notes
    \item \texttt{cultural\_weight}: Importance score (0.0-1.0) based on expert assessment
    \item \texttt{thematic\_category}: Primary cultural theme (e.g., ``Reciprocity'', ``Patience'')
\end{itemize}

\paragraph{CulturalConcept Nodes} (\texttt{:CulturalConcept}, 847 instances) represent abstract cultural themes, values, and knowledge domains:
\begin{itemize}
    \item \texttt{concept\_name}: Canonical name (e.g., ``Reciprocity'', ``NgwatioSystem'')
    \item \texttt{definition}: Formal definition of the cultural concept
    \item \texttt{cultural\_significance}: Anthropological context and importance
    \item \texttt{business\_relevance}: Contemporary economic interpretations
    \item \texttt{hierarchy\_level}: Depth in concept taxonomy (0 = root, 1 = theme, 2+ = sub-concepts)
\end{itemize}

\paragraph{UsageContext Nodes} (\texttt{:UsageContext}, 31 instances) encode appropriate situations for proverb use:
\begin{itemize}
    \item \texttt{context\_name}: Context identifier (e.g., ``WealthDistributionCeremony'')
    \item \texttt{description}: Detailed scenario description
    \item \texttt{participants}: Typical social roles involved
    \item \texttt{appropriateness}: When this context is culturally appropriate
\end{itemize}

\paragraph{MoralLesson Nodes} (\texttt{:MoralLesson}, 43 instances) capture ethical principles:
\begin{itemize}
    \item \texttt{lesson\_name}: Principle name (e.g., ``CommunityOverIndividual'')
    \item \texttt{teaching}: Explicit moral instruction
    \item \texttt{examples}: Illustrative scenarios
\end{itemize}

\subsubsection{Relationship Types}
\label{subsubsec:relationship_types}

Six relationship types link nodes, forming the semantic network:

\begin{description}
    \item[\texttt{:EXPRESSES\_CONCEPT}] (Proverb $\rightarrow$ CulturalConcept): Indicates a proverb expresses a cultural theme. Properties: \texttt{salience} (0.0-1.0 indicating centrality of concept to proverb meaning).
    
    \item[\texttt{:TEACHES\_LESSON}] (Proverb $\rightarrow$ MoralLesson): Links proverbs to moral teachings. Properties: \texttt{directness} (explicit vs. implicit teaching).
    
    \item[\texttt{:USED\_IN}] (Proverb $\rightarrow$ UsageContext): Specifies appropriate usage contexts. Properties: \texttt{appropriateness\_score} (cultural fit rating).
    
    \item[\texttt{:RELATES\_TO}] (CulturalConcept $\leftrightarrow$ CulturalConcept): Associative connections between concepts (e.g., Reciprocity relates to CommunityValues). Properties: \texttt{strength} (relatedness score).
    
    \item[\texttt{:SUBSUMES}] (CulturalConcept $\rightarrow$ CulturalConcept): Hierarchical parent-child relationships (e.g., EconomicConcepts subsumes Reciprocity). Properties: \texttt{is\_defining} (whether child is definitional for parent).
    
    \item[\texttt{:REFERENCES}] (Proverb $\rightarrow$ CulturalEntity): Links to concrete cultural referents (animals, plants, historical figures). Properties: \texttt{metaphorical} (boolean indicating metaphorical vs. literal reference).
\end{description}

\subsubsection{Constraints and Indexes}
\label{subsubsec:constraints_indexes}

To ensure data integrity and query performance, the schema defines:

\paragraph{Uniqueness Constraints}
\begin{verbatim}
CREATE CONSTRAINT proverb_id IF NOT EXISTS
FOR (p:Proverb) REQUIRE p.proverb_id IS UNIQUE;

CREATE CONSTRAINT concept_name IF NOT EXISTS
FOR (c:CulturalConcept) REQUIRE c.concept_name IS UNIQUE;
\end{verbatim}

\paragraph{Property Indexes} for efficient lookups:
\begin{verbatim}
CREATE INDEX proverb_weight IF NOT EXISTS
FOR (p:Proverb) ON (p.cultural_weight);

CREATE INDEX concept_hierarchy IF NOT EXISTS
FOR (c:CulturalConcept) ON (c.hierarchy_level);
\end{verbatim}

\paragraph{Full-Text Indexes} enabling semantic search:
\begin{verbatim}
CREATE FULLTEXT INDEX proverb_text IF NOT EXISTS
FOR (p:Proverb) ON EACH [p.kikuyu_text, 
                          p.expert_translation,
                          p.expert_cultural_meaning];

CREATE FULLTEXT INDEX concept_definitions IF NOT EXISTS
FOR (c:CulturalConcept) ON EACH [c.concept_name,
                                   c.definition,
                                   c.cultural_significance];
\end{verbatim}

\subsection{Ontology Population Pipeline}
\label{subsec:ontology_population}

Populating the knowledge graph from raw proverb data and expert annotations involved a multi-stage ETL (Extract-Transform-Load) pipeline implemented in \texttt{scripts/ontology\_builder.py}:

\begin{enumerate}
    \item \textbf{Data Extraction}: Parsing Ireri's proverb collection from CSV format, extracting Kikuyu text, expert translations, cultural meanings, and business relevance notes.
    
    \item \textbf{Concept Identification}: Using keyword matching and LLM-assisted extraction (\texttt{extract\_ontology\_concepts\_with\_llm.py}) to identify cultural concepts referenced in each proverb. Example: for proverb MW\_002 (``Andu ni indo''), concepts ``Community'', ``SocialCapital'', ``ValueHierarchy'' were extracted.
    
    \item \textbf{Relationship Inference}: Creating \texttt{:EXPRESSES\_CONCEPT} edges with salience scores computed from term frequency in cultural meaning annotations. High-salience concepts (score $>$ 0.7) become primary themes for retrieval.
    
    \item \textbf{Weight Assignment}: Computing \texttt{cultural\_weight} for proverbs through expert surveys (3 Kikuyu elders rated importance on 1-5 scale, normalized to 0.0-1.0). High-weight proverbs (e.g., MW\_002 at 0.92) are prioritized during retrieval.
    
    \item \textbf{Validation}: Running \texttt{scripts/ontology\_validator.py} to check: (a) all proverbs link to at least one concept, (b) concept hierarchies are acyclic, (c) no orphaned nodes exist.
\end{enumerate}

The final populated graph contains 947 nodes (100 Proverbs + 847 CulturalConcepts + supplementary nodes) and 1,247 edges, achieving 99.8\% schema compliance (2 minor anomalies in concept hierarchy depth, manually corrected).

\section{Hybrid Retrieval Engine}
\label{sec:retrieval_engine}

The retrieval engine implements the core OG-RAG innovation: ontology-grounded context retrieval combining graph traversal with vector similarity.

\subsection{Retrieval Strategy}
\label{subsec:retrieval_strategy}

For a given Kikuyu proverb input, the retriever executes a three-phase strategy:

\subsubsection{Phase 1: Concept Extraction}
\label{subsubsec:concept_extraction}

Input text is analyzed to identify cultural concepts using two methods:

\paragraph{Keyword Matching} A curated dictionary maps lexical patterns to concepts:
\begin{verbatim}
CONCEPT_KEYWORDS = {
    'wealth': ['utonga', 'mbeca', 'indo', 'gitonga'],
    'reciprocity': ['ngwatio', 'ũtugi', 'reciprocal'],
    'patience': ['kĩrĩa', 'eterera', 'wait'],
    ...
}
\end{verbatim}

Matching against proverb text yields candidate concepts (e.g., ``wealth'', ``patience'').

\paragraph{Semantic Expansion} Matched concepts are expanded to include hierarchically related concepts via graph traversal:
\begin{verbatim}
MATCH (c:CulturalConcept {concept_name: $matched_concept})
      -[:RELATES_TO|SUBSUMES*1..2]-(related:CulturalConcept)
RETURN DISTINCT related.concept_name
\end{verbatim}

This expands ``wealth'' to include ``WealthAcquisition'', ``MaterialProsperity'', ``EconomicWisdom'', etc., capturing broader semantic neighborhoods.

\subsubsection{Phase 2: Graph-Based Retrieval}
\label{subsubsec:graph_retrieval}

Using extracted concepts, Cypher queries retrieve proverbs expressing similar cultural themes:

\begin{verbatim}
MATCH (c:CulturalConcept)<-[r:EXPRESSES_CONCEPT]-(p:Proverb)
WHERE c.concept_name IN $extracted_concepts
  AND NOT p.proverb_id = $input_proverb_id
WITH p, 
     collect(DISTINCT c.concept_name) as matched_concepts,
     sum(r.salience) as concept_score
RETURN p, matched_concepts, concept_score
ORDER BY concept_score DESC
LIMIT $k
\end{verbatim}

This retrieves the top-$k$ proverbs sharing the most conceptual overlap (weighted by salience) with the input.

\subsubsection{Phase 3: Hybrid Scoring and Re-ranking}
\label{subsubsec:hybrid_scoring}

Graph-retrieved proverbs are re-ranked using a hybrid score combining three signals:

\begin{equation}
\text{HybridScore}(p) = w_c \cdot S_{\text{concept}}(p) + w_w \cdot S_{\text{weight}}(p) + w_l \cdot S_{\text{lexical}}(p)
\end{equation}

where:
\begin{itemize}
    \item $S_{\text{concept}}(p)$: Normalized concept overlap score from Phase 2
    \item $S_{\text{weight}}(p)$: Cultural weight of proverb $p$ (expert importance rating)
    \item $S_{\text{lexical}}(p)$: Jaccard similarity between Kikuyu text of $p$ and input proverb
    \item $w_c = 0.5$, $w_w = 0.3$, $w_l = 0.2$: Empirically tuned weights
\end{itemize}

This hybrid approach balances semantic similarity (concept overlap) with cultural significance (expert weights) and surface-level textual similarity (lexical matching), addressing the limitations of pure vector similarity discussed in Section~\ref{subsec:tech_challenge}.

\subsection{Implementation: GraphRetriever Class}
\label{subsec:graph_retriever_impl}

The retrieval logic is encapsulated in \texttt{src/og-rag-system/graph\_retriever.py}:

\begin{verbatim}
class GraphRetriever:
    def __init__(self, uri, username, password, database='neo4j'):
        self.driver = GraphDatabase.driver(uri, 
                                           auth=(username, password))
        self.database = database
    
    def extract_concepts(self, text: str) -> List[str]:
        """Extract cultural concepts from text."""
        # Keyword matching + semantic expansion
        ...
    
    def retrieve_by_concepts(self, concepts: List[str], 
                            k: int = 5) -> List[RetrievedProverb]:
        """Retrieve top-k proverbs via graph query."""
        # Cypher query execution
        ...
    
    def hybrid_retrieve(self, kikuyu_text: str, 
                       k: int = 5) -> List[RetrievedProverb]:
        """Full hybrid retrieval pipeline."""
        concepts = self.extract_concepts(kikuyu_text)
        candidates = self.retrieve_by_concepts(concepts, k=k*3)
        reranked = self._hybrid_score(candidates, kikuyu_text)
        return reranked[:k]
\end{verbatim}

The \texttt{RetrievedProverb} dataclass bundles retrieved proverb data with metadata:
\begin{verbatim}
@dataclass
class RetrievedProverb:
    proverb_id: str
    kikuyu_text: str
    expert_translation: str
    expert_cultural_meaning: str
    cultural_weight: float
    similarity_score: float
    matched_concepts: List[str]
    retrieval_method: str  # 'graph', 'hybrid', 'lexical'
\end{verbatim}

\subsection{Retrieval Performance}
\label{subsec:retrieval_performance}

Retrieval latency was profiled across the 100-proverb corpus:
\begin{itemize}
    \item \textbf{Concept Extraction}: 0.12s average (keyword matching dominates)
    \item \textbf{Graph Query}: 0.68s average (Cypher execution + network latency to AuraDB)
    \item \textbf{Hybrid Scoring}: 0.22s average (Python-side computation)
    \item \textbf{Total Retrieval}: 1.02s average per proverb
\end{itemize}

For production deployment, caching frequently retrieved subgraphs and precomputing concept embeddings could reduce latency by an estimated 40-60\%.

\section{Context Builder}
\label{sec:context_builder}

The Context Builder transforms graph-retrieved knowledge into natural language prompts optimized for LLM consumption.

\subsection{Context Serialization Strategy}
\label{subsec:context_serialization}

Retrieved graph subgraphs (nodes + relationships) are serialized into structured text following a five-section template:

\begin{enumerate}
    \item \textbf{Cultural Themes}: Lists extracted concepts with definitions
    \item \textbf{Related Proverbs}: Top-3 similar proverbs with expert translations and cultural meanings
    \item \textbf{Usage Contexts}: Appropriate scenarios for proverb application
    \item \textbf{Moral Lessons}: Ethical principles the proverb teaches
    \item \textbf{Metaphorical Elements}: Cultural entities referenced and their symbolic meanings
\end{enumerate}

Example serialization for input ``Andu ni indo'' (People are wealth):

\begin{verbatim}
CULTURAL CONTEXT (from Kikuyu Proverb Ontology):

CULTURAL THEMES:
- Community: Collective identity and mutual support systems
- Social Capital: Human relationships as wealth
- Value Hierarchy: Kikuyu priorities of people over possessions

RELATED PROVERBS:
1. "Ciringi imurikagwa na ingi" (Money is made by money)
   Cultural Meaning: Wealth multiplication requires initial capital
   
2. "Gutiri mbura itari gitonga kiayo" (No rain without its wealth)
   Cultural Meaning: Every challenge brings opportunities

USAGE CONTEXTS:
- Wealth distribution ceremonies (emphasizing communal sharing)
- Advice to young people (teaching value priorities)
- Business partnerships (prioritizing trust over capital)

MORAL LESSONS:
- People constitute ultimate wealth, not material possessions
- Community relationships generate sustainable prosperity
- Social obligations precede individual accumulation
\end{verbatim}

This structured format enables LLMs to parse specific knowledge types (themes vs. examples vs. contexts) rather than processing an undifferentiated text blob.

\subsection{Implementation: ContextBuilder Class}
\label{subsec:context_builder_impl}

Implemented in \texttt{src/og-rag-system/context\_builder.py}:

\begin{verbatim}
class ContextBuilder:
    def __init__(self, graph_retriever: GraphRetriever):
        self.retriever = graph_retriever
    
    def build_cultural_context(
        self, 
        kikuyu_text: str, 
        k_proverbs: int = 3
    ) -> CulturalContext:
        """Build structured context from graph retrieval."""
        
        # Retrieve similar proverbs
        retrieved = self.retriever.hybrid_retrieve(
            kikuyu_text, k=k_proverbs
        )
        
        # Extract unique concepts
        all_concepts = set()
        for p in retrieved:
            all_concepts.update(p.matched_concepts)
        
        # Build sections
        context = CulturalContext(
            themes=self._format_themes(all_concepts),
            related_proverbs=self._format_proverbs(retrieved),
            contexts=self._format_usage_contexts(all_concepts),
            lessons=self._format_moral_lessons(all_concepts)
        )
        
        return context
    
    def to_prompt_string(self, context: CulturalContext) -> str:
        """Serialize context to LLM prompt format."""
        # Template-based string generation
        ...
\end{verbatim}

The \texttt{CulturalContext} dataclass stores structured sections before final serialization, enabling A/B testing of different prompt formats without modifying retrieval logic.

\section{LLM Integration Layer}
\label{sec:llm_integration}

The LLM Integration Layer interfaces with OpenAI and Google Gemini APIs to generate translations guided by ontology-derived context.

\subsection{Translation Modes}
\label{subsec:translation_modes}

The system implements three translation modes for comparative evaluation:

\subsubsection{Raw LLM Mode}
\label{subsubsec:raw_llm}

Baseline approach using zero-shot prompting without external context:

\begin{verbatim}
PROMPT:
You are an expert translator specializing in African languages.
Translate the following Kikuyu proverb to English, preserving
its cultural meaning and wisdom.

Kikuyu: "{kikuyu_text}"

Provide:
1. English translation
2. Brief cultural explanation
\end{verbatim}

This tests the LLM's intrinsic multilingual and cultural knowledge.

\subsubsection{Traditional RAG Mode}
\label{subsubsec:traditional_rag}

Standard RAG approach retrieving similar proverbs via vector similarity (Sentence-BERT embeddings) without ontological structure:

\begin{verbatim}
PROMPT:
You are an expert translator. Use these example Kikuyu proverbs
to guide your translation:

EXAMPLES:
{vector_retrieved_examples}

TRANSLATE:
Kikuyu: "{kikuyu_text}"
\end{verbatim}

This tests whether unstructured retrieval provides sufficient cultural context.

\subsubsection{OG-RAG Mode}
\label{subsubsec:ograg_mode}

Full ontology-grounded approach with structured cultural context:

\begin{verbatim}
PROMPT:
You are an expert in Kikuyu culture specializing in proverb
translation. Use the following ontology-derived cultural 
knowledge to produce a culturally faithful translation.

{structured_cultural_context}

PROVERB TO TRANSLATE:
Kikuyu: "{kikuyu_text}"
Literal gloss: "{literal_translation}"

REQUIREMENTS:
1. Preserve cultural worldview (not just literal meaning)
2. Maintain metaphorical imagery where appropriate
3. Adapt culturally-specific references for English audiences
4. Ensure translation would be appropriate in usage contexts
   listed above

Provide:
1. Culturally faithful English translation
2. Brief explanation of cultural adaptation choices
\end{verbatim}

This represents the core OG-RAG contribution.

\subsection{Implementation: OGRAGTranslator Class}
\label{subsec:ograg_translator_impl}

The end-to-end translation pipeline is implemented in \texttt{src/og-rag-system/ograg\_translator.py}:

\begin{verbatim}
class OGRAGTranslator:
    def __init__(self, openai_api_key: str, 
                 model: str = 'gpt-4-turbo',
                 temperature: float = 0.3):
        self.client = OpenAI(api_key=openai_api_key)
        self.model = model
        self.temperature = temperature
        self.graph_retriever = GraphRetriever(...)
        self.context_builder = ContextBuilder(self.graph_retriever)
    
    def translate_raw(self, kikuyu_text: str) -> TranslationResult:
        """Raw LLM translation (baseline)."""
        prompt = self._build_raw_prompt(kikuyu_text)
        response = self.client.chat.completions.create(
            model=self.model,
            messages=[{"role": "user", "content": prompt}],
            temperature=self.temperature,
            max_tokens=500
        )
        return self._parse_response(response, method='raw')
    
    def translate_traditional_rag(self, kikuyu_text: str) 
                                 -> TranslationResult:
        """Traditional RAG with vector similarity."""
        examples = self._retrieve_vector_examples(kikuyu_text)
        prompt = self._build_rag_prompt(kikuyu_text, examples)
        response = self.client.chat.completions.create(...)
        return self._parse_response(response, 
                                    method='traditional_rag')
    
    def translate_ograg(self, kikuyu_text: str) -> TranslationResult:
        """OG-RAG with ontology-grounded context."""
        cultural_context = self.context_builder.build_cultural_context(
            kikuyu_text, k_proverbs=3
        )
        prompt = self._build_ograg_prompt(kikuyu_text, 
                                          cultural_context)
        response = self.client.chat.completions.create(...)
        return self._parse_response(response, method='ograg')
\end{verbatim}

The \texttt{TranslationResult} dataclass captures comprehensive metadata for analysis:
\begin{verbatim}
@dataclass
class TranslationResult:
    proverb_id: str
    kikuyu_text: str
    translation: str
    explanation: Optional[str]
    method: str  # 'raw', 'traditional_rag', 'ograg'
    retrieved_proverbs: Optional[List[RetrievedProverb]]
    concepts_used: Optional[List[str]]
    prompt_tokens: int
    completion_tokens: int
    total_tokens: int
    timestamp: str
\end{verbatim}

\subsection{API Management and Cost Optimization}
\label{subsec:api_management}

OpenAI API costs were managed through:
\begin{itemize}
    \item \textbf{Model Selection}: Using GPT-4-turbo (\$0.01/1K input tokens, \$0.03/1K output) for OG-RAG vs. GPT-3.5-turbo (\$0.0005/1K) for preliminary testing
    \item \textbf{Response Caching}: Storing translations in \texttt{data/results/} to avoid redundant API calls during analysis iterations
    \item \textbf{Token Budgeting}: Limiting max\_tokens=500 for translations (empirically sufficient for proverb translations + explanations)
    \item \textbf{Batch Processing}: Processing 100 proverbs in batches of 10 with exponential backoff for rate limiting
\end{itemize}

Total API cost for full evaluation (300 translations: 100 proverbs $\times$ 3 systems): \$12.47 (within research budget).

\section{Baseline Systems}
\label{sec:baseline_systems}

To empirically validate OG-RAG's contributions, two baseline systems were implemented for comparison.

\subsection{Raw GPT-4 Baseline}
\label{subsec:raw_baseline}

Implemented in \texttt{src/evaluation/baseline\_translation\_system.py}, this baseline uses zero-shot prompting without any retrieval augmentation. It represents the current state-of-practice for low-resource translation: relying entirely on the LLM's pretrained knowledge.

Performance characteristics:
\begin{itemize}
    \item \textbf{Latency}: 1.8s average (single API call, no retrieval overhead)
    \item \textbf{Token Usage}: 127 prompt tokens, 89 completion tokens average
    \item \textbf{Cost}: \$0.0037 per translation
\end{itemize}

\subsection{Traditional RAG Baseline}
\label{subsec:traditional_rag_baseline}

This baseline implements standard RAG using Sentence-BERT embeddings to retrieve similar proverbs:

\begin{enumerate}
    \item \textbf{Embedding Generation}: All 100 proverbs (Kikuyu + English) embedded using \texttt{sentence-transformers/paraphrase-multilingual-mpnet-base-v2}
    \item \textbf{Vector Similarity Search}: Cosine similarity between input embedding and corpus embeddings
    \item \textbf{Context Injection}: Top-3 similar proverbs (by vector distance) included in prompt as examples
\end{enumerate}

Performance characteristics:
\begin{itemize}
    \item \textbf{Latency}: 2.4s average (0.5s embedding + 0.1s similarity search + 1.8s LLM generation)
    \item \textbf{Token Usage}: 312 prompt tokens, 94 completion tokens average (higher due to example context)
    \item \textbf{Cost}: \$0.0061 per translation
\end{itemize}

Key limitation: Vector similarity captures surface-level semantic overlap but misses structured cultural relationships (e.g., that reciprocity relates to community values, not merely via keyword co-occurrence).

\section{Evaluation Pipeline}
\label{sec:eval_pipeline}

The evaluation pipeline orchestrates translation generation, metric computation, and statistical analysis.

\subsection{Comparative Evaluation Framework}
\label{subsec:comparative_framework}

Implemented in \texttt{src/evaluation/comparative\_pipeline.py}, the framework:

\begin{enumerate}
    \item Loads the 100-proverb evaluation corpus
    \item Generates translations using all three systems (Raw, Traditional RAG, OG-RAG)
    \item Anonymizes and randomizes translations for blind evaluation
    \item Computes automated metrics (cultural authenticity, translation fidelity, overall quality)
    \item Performs statistical significance tests (paired t-tests, effect sizes)
    \item Generates visualization outputs
\end{enumerate}

\subsection{Metric Computation}
\label{subsec:metric_computation}

The \texttt{src/evaluation/cultural\_metrics.py} module implements the two-dimensional evaluation framework from Section~\ref{subsec:eval_metrics}:

\begin{verbatim}
class CulturalMetrics:
    def compute_cultural_authenticity(
        self, 
        translation: str,
        expert_cultural_meaning: str,
        cultural_themes: List[str]
    ) -> float:
        """
        Compute cultural authenticity score (0.0-1.0).
        
        Components:
        - Thematic accuracy (40%): Theme preservation
        - Metaphorical preservation (30%): Figurative language
        - Contextual appropriateness (20%): Usage contexts
        - Value alignment (10%): Cultural values
        """
        # LLM-assisted scoring against expert annotations
        ...
    
    def compute_translation_fidelity(
        self,
        translation: str,
        expert_translation: str,
        kikuyu_text: str
    ) -> float:
        """
        Compute translation fidelity score (0.0-1.0).
        
        Components:
        - Semantic accuracy (50%): Meaning preservation
        - Fluency (30%): English naturalness
        - Completeness (20%): Information coverage
        """
        # Combination of BLEU, semantic similarity, fluency
        ...
    
    def compute_overall_quality(
        self,
        cultural_authenticity: float,
        translation_fidelity: float
    ) -> float:
        """Weighted composite: 60% cultural + 40% fidelity."""
        return 0.6 * cultural_authenticity + 0.4 * translation_fidelity
\end{verbatim}

\subsection{Statistical Analysis Module}
\label{subsec:statistical_analysis}

The \texttt{src/evaluation/statistical\_analysis.py} module implements the statistical testing framework:

\begin{verbatim}
from scipy import stats

def paired_t_test(scores_a: List[float], 
                 scores_b: List[float]) -> Tuple[float, float]:
    """
    Perform paired t-test for system comparison.
    
    Returns:
        (t_statistic, p_value)
    """
    t_stat, p_val = stats.ttest_rel(scores_a, scores_b)
    return t_stat, p_val

def cohens_d(scores_a: List[float], 
            scores_b: List[float]) -> float:
    """
    Compute Cohen's d effect size.
    
    Interpretation:
    - Small: d ~ 0.2
    - Medium: d ~ 0.5
    - Large: d ~ 0.8
    """
    mean_diff = np.mean(scores_a) - np.mean(scores_b)
    pooled_std = np.sqrt((np.var(scores_a) + np.var(scores_b)) / 2)
    return mean_diff / pooled_std
\end{verbatim}

\section{Implementation Challenges and Solutions}
\label{sec:implementation_challenges}

Several technical challenges emerged during implementation:

\subsection{Neo4j AuraDB Cloud Latency}
\label{subsec:auradb_latency}

\textbf{Challenge}: Cloud-hosted Neo4j AuraDB introduced 300-500ms network latency per query, significantly impacting retrieval performance.

\textbf{Solution}: Implemented query result caching using Python's \texttt{functools.lru\_cache}, reducing repeat queries for identical concept sets. For production deployment, migrating to a self-hosted Neo4j instance could eliminate network latency.

\subsection{Concept Extraction Ambiguity}
\label{subsec:concept_ambiguity}

\textbf{Challenge}: Keyword-based concept extraction produced false positives (e.g., ``muti'' matching both ``tree'' and ``medicine'' concepts when only one was intended).

\textbf{Solution}: Implemented context-aware disambiguation using surrounding words and synonym lists. For example, ``muti wa ũganga'' (medicine tree) correctly maps to ``TraditionalMedicine'' not ``ForestResources''.

\subsection{LLM Response Parsing}
\label{subsec:response_parsing}

\textbf{Challenge}: LLM outputs occasionally deviated from expected format (translation + explanation), complicating automated parsing.

\textbf{Solution}: Employed structured prompting with explicit section markers (``TRANSLATION:'', ``EXPLANATION:'') and regex-based post-processing to extract components. Fallback handling logs unparseable responses for manual review (occurred in 3\% of cases).

\subsection{Cultural Weight Calibration}
\label{subsec:weight_calibration}

\textbf{Challenge}: Expert-assigned cultural weights showed high inter-rater variance (Krippendorff's $\alpha = 0.61$, below ideal threshold).

\textbf{Solution}: Conducted adjudication sessions where experts discussed discrepancies, achieving consensus weights. Final weights showed improved consistency ($\alpha = 0.78$).

\section{Code Organization and Repository Structure}
\label{sec:code_organization}

The codebase is organized for modularity and reproducibility:

\begin{verbatim}
opit-rai9001/
├── src/
│   ├── ontology/
│   │   ├── enhanced_neo4j_schema.py      # Schema definition
│   │   ├── cultural_weights.py           # Weight assignments
│   │   └── kikuyu_proverb_ontology.py    # Ontology utils
│   ├── og-rag-system/
│   │   ├── graph_retriever.py            # Hybrid retrieval
│   │   ├── context_builder.py            # Context serialization
│   │   └── ograg_translator.py           # Translation pipeline
│   ├── evaluation/
│   │   ├── cultural_metrics.py           # Metric computation
│   │   ├── statistical_analysis.py       # Statistical tests
│   │   ├── baseline_translation_system.py # Baselines
│   │   └── comparative_pipeline.py       # Evaluation orchestration
│   └── data_loading/
│       └── proverb_loader.py             # Corpus ingestion
├── scripts/
│   ├── ontology_builder.py               # KG population
│   ├── ontology_validator.py             # Validation checks
│   └── run_evaluation.py                 # Main evaluation script
├── data/
│   ├── proverbs/                         # Source corpus
│   ├── ontology/                         # Ontology artifacts
│   ├── results/                          # Translation outputs
│   └── evaluation/                       # Metrics and analysis
├── docs/thesis/                          # Thesis LaTeX files
└── requirements.txt                      # Python dependencies
\end{verbatim}

All code is version-controlled in GitHub (\url{https://github.com/ndethi/opit-rai9001}), with tagged releases corresponding to experimental runs reported in Chapter~\ref{chap:evaluation}.

\section{Summary}
\label{sec:design_summary}

This chapter detailed the technical architecture and implementation of the thiLLMo OG-RAG system, spanning knowledge graph infrastructure, hybrid retrieval mechanisms, context building, LLM integration, and evaluation pipelines. The modular design enables systematic comparison of OG-RAG against baseline approaches (Raw LLM, Traditional RAG) while maintaining transparency and reproducibility.

Key architectural contributions include: (1) a multi-layered Neo4j schema encoding 847 cultural concepts with rich relational semantics, (2) a hybrid retrieval strategy combining graph traversal, cultural weight signals, and lexical similarity to achieve ontology-grounded context selection, (3) a structured context serialization approach transforming graph data into LLM-optimized prompts, and (4) a comprehensive evaluation framework operationalizing cultural fidelity measurement.

The implementation faced challenges typical of cultural NLP research—ambiguous concept extraction, expert knowledge elicitation, LLM output variability—addressed through context-aware disambiguation, structured adjudication processes, and robust parsing mechanisms. The resulting system provides a replicable platform for investigating how structured cultural knowledge enhances LLM-based translation quality.

The next chapter (Chapter~\ref{chap:evaluation}) presents the empirical evaluation of this system, demonstrating through rigorous quantitative and qualitative analysis that ontology-grounded retrieval significantly improves both cultural authenticity and translation fidelity compared to baseline approaches.
